\UnnumberedChapter{KẾT LUẬN}

Sau quá trình tìm hiểu và triển khai, nhóm chúng em đã xây dựng được một AI Vision Service hoàn chỉnh cho hệ thống Smart Campus, từ khâu tiếp nhận ảnh đầu vào, xử lý bằng mô hình nhận diện, chuẩn hóa kết quả trả về cho đến đóng gói và mô tả API để các dịch vụ khác có thể tích hợp thuận tiện. Đề tài cho thấy việc tách riêng một dịch vụ AI ra khỏi các thành phần nghiệp vụ khác là một hướng tiếp cận hợp lý, giúp hệ thống dễ mở rộng, dễ kiểm thử và linh hoạt hơn trong triển khai thực tế.

Trong quá trình thực hiện, nhóm đã hoàn thiện các phần quan trọng như cấu trúc service, API contract, OpenAPI specification, Postman collection, smoke test và tài liệu báo cáo. Đặc biệt, cơ chế fallback mock AI giúp hệ thống vẫn có thể hoạt động trong những trường hợp mô hình thật chưa sẵn sàng, từ đó tăng tính ổn định cho toàn bộ nền tảng. Bên cạnh đó, việc tổ chức tài liệu theo template cũng giúp quá trình trình bày, tổng hợp và bảo trì nội dung được thuận lợi hơn.

Từ đề tài này, nhóm chúng em rút ra được nhiều kinh nghiệm về cách thiết kế một service AI trong kiến trúc microservice, cách mô tả rõ ràng hợp đồng dữ liệu giữa các thành phần và cách chuẩn hóa quy trình tích hợp giữa camera, AI và các dịch vụ nghiệp vụ. Đây là những kiến thức có giá trị không chỉ cho bài tập lớn hiện tại mà còn cho các hướng phát triển tiếp theo trong các hệ thống giám sát thông minh.

Trong thời gian tới, nếu tiếp tục mở rộng đề tài, nhóm có thể bổ sung cơ chế lưu trữ sự kiện lâu dài, cải tiến chất lượng mô hình nhận diện, tối ưu hiệu năng xử lý ảnh và hoàn thiện thêm các kênh đồng bộ dữ liệu với những dịch vụ khác của Smart Campus. Nhóm chúng em xin chân thành cảm ơn cô giáo hướng dẫn và các thầy cô trong khoa đã hỗ trợ để báo cáo này được hoàn thành.
